\DocumentMetadata{
  pdfversion=2.0,
  pdfstandard=ua-2,
  lang=en-US,
  tagging=on,
  tagging-setup={math/setup=mathml-SE,tabsorder=structure}
}
\documentclass[11pt,letterpaper]{article}
\usepackage[margin=0.68in,includefoot]{geometry}
\usepackage{newcomputermodern}
\usepackage{amsmath,amscd,mathtools}
\providecommand{\square}{\mdlgwhtsquare}
\usepackage[hidelinks]{hyperref}
\hypersetup{
  pdftitle={Probability PhD Exam, August 2023},
  pdfauthor={University of Florida Department of Mathematics},
  pdfsubject={Graduate mathematics examination},
  pdfdisplaydoctitle=true,
  bookmarksopen=true
}
\setcounter{secnumdepth}{0}
\setlength{\parindent}{0pt}
\setlength{\parskip}{0.28em}
\setlength{\emergencystretch}{3em}
\setlength{\leftmargini}{2.15em}
\setlength{\leftmarginii}{2.65em}
\setlength{\leftmarginiii}{3em}
\renewcommand{\labelenumi}{\textbf{\theenumi.}}
\pagestyle{empty}
\raggedbottom
\allowdisplaybreaks
\newenvironment{examproblems}[1]{%
  \begin{enumerate}%
  \setcounter{enumi}{#1}%
  \setlength{\topsep}{0.35em}%
  \setlength{\partopsep}{0pt}%
  \setlength{\itemsep}{0.48em}%
  \setlength{\parsep}{0.18em}%
}{\end{enumerate}}
\newenvironment{examsubparts}{%
  \begin{enumerate}%
  \setlength{\topsep}{0.18em}%
  \setlength{\partopsep}{0pt}%
  \setlength{\itemsep}{0.16em}%
  \setlength{\parsep}{0.1em}%
}{\end{enumerate}}
\newenvironment{examinstructions}{%
  \par\begingroup\itshape
}{%
  \par\endgroup\vspace{0.18em}
}
\makeatletter
\renewcommand\section{\@startsection{section}{1}{0pt}%
  {-1.25ex plus -.25ex minus -.1ex}%
  {0.55ex plus .1ex}%
  {\normalfont\large\bfseries}}
\renewcommand{\@maketitle}{%
  \newpage\null\vskip -1.2em%
  {\footnotesize\bfseries
    \href{https://www.ufl.edu/}{University of Florida}, \href{https://math.ufl.edu/}{Department of Mathematics}\par}%
  \vskip 0.18em%
  \tagstructbegin{tag=Title}%
    {\LARGE\bfseries\href{https://gma.math.ufl.edu/exams/probability-phd/}{Probability PhD Exam}, August 2023\par}%
  \tagstructend%
  \vskip 0.35em%
  \tagmcbegin{artifact}%
    \hrule height 1pt%
  \tagmcend%
  \vskip 0.55em%
}
\makeatother
\title{Probability PhD Exam, August 2023}
\author{}
\date{}
\begin{document}
\maketitle
\thispagestyle{empty}
\begin{examinstructions}
Carefully justify your answers.
\end{examinstructions}
\begin{examproblems}{0}
\item
Let \(X,Y:(\Omega,\mathcal{F},\mathbb{P})\to(\mathbb{R},\mathcal{B}(\mathbb{R}))\) be random variables. Prove that 
\[
\mathbb{P}(X+Y\ge 0)\le \mathbb{P}(X\ge 0)+\mathbb{P}(Y\ge 0).
\]
\item
This problem has two parts.
\begin{examsubparts}
\item[\textnormal{1.}]
Give the mathematical definition of:
\begin{examsubparts}
\item[\textnormal{(a)}]
convergence almost sure,
\item[\textnormal{(b)}]
convergence in probability,
\item[\textnormal{(c)}]
convergence in \(L^1\),
\item[\textnormal{(d)}]
convergence in distribution.
\end{examsubparts}
\item[\textnormal{2.}]
For \(n\ge 1\), let \(X_n\) be uniformly distributed on \(\{1,\ldots,n\}\), that is 
\[
\mathbb{P}(X_n=k)=\frac{1}{n},\qquad k\in\{1,\ldots,n\}.
\]
 Prove that \(\left\{\frac{X_n}{n}\right\}_{n\ge 1}\) converges to~\(U\) in distribution, where~\(U\) is uniform on \([0,1]\).
\end{examsubparts}
\item
Let \(\{X_i\}_{i\ge 1}\) be a sequence of i.i.d. continuous random variables having a uniform distribution on \([0,1]\). Define, for \(n\ge 1\), 
\[
Y_n=\max(X_1,\ldots,X_n).
\]
\begin{examsubparts}
\item[\textnormal{1.}]
Find the cumulative distribution function (CDF) of \(Y_n\), \(n\ge 1\).
\item[\textnormal{2.}]
Compute \(\mathbb{E}[Y_n]\) and \(\operatorname{Var}(Y_n)\).
\item[\textnormal{3.}]
Prove that \(\{Y_n\}\) converges to~\(1\) in \(L^1\). What about almost sure convergence?
\end{examsubparts}
\item
Let \(\{X_i\}_{i\ge 1}\) be a sequence of i.i.d. random variables with a Poisson distribution of parameter~\(\lambda\).
\begin{examsubparts}
\item[\textnormal{1.}]
Let \(n\ge 1\). Without justification, what is the distribution of \(X_1+\cdots+X_n\)?
\item[\textnormal{2.}]
State the weak and strong law of large numbers for i.i.d random variables.
\item[\textnormal{3.}]
Let \(\phi:\mathbb{R}\to\mathbb{R}\) be an arbitrary bounded continuous function. Prove that 
\[
\lim_{n\to+\infty}\sum_{k=0}^{+\infty}e^{-n\lambda}\frac{(n\lambda)^k}{k!}\phi\left(\frac{k}{n}\right)=\phi(\lambda).
\]
\end{examsubparts}
\item
This problem has three parts.
\begin{examsubparts}
\item[\textnormal{1.}]
Give the mathematical definition of standard Brownian motion.
\item[\textnormal{2.}]
Let \(\{B_t\}\) be a standard Brownian motion. Prove that \(\left\{\frac{B_n}{n}\right\}\) converges to~\(0\) almost surely as \(n\to+\infty\) (\(n\in\mathbb{N}\)).
\item[\textnormal{3.}]
Find \(\mathbb{P}(B_2\ge 0,B_1\le 0)\).
\end{examsubparts}
\item
Let \(\{M_n\}_{n\ge 0}\) be a sequence of integrable random variables adapted to a filtration \(\{\mathcal{F}_n\}\). Assume that for each bounded stopping time~\(T\), \(\mathbb{E}[M_T]=\mathbb{E}[M_0]\). Show that \(\{M_n\}_{n\ge 0}\) is a martingale.
\item
Let \(\{B_t\}\) be a standard Brownian motion. Define, for \(x\in\mathbb{R}\), 
\[
T_x=\min\{t\ge 0:B_t=x\}.
\]
\begin{examsubparts}
\item[\textnormal{1.}]
Prove that for all \(x\in\mathbb{R}\), \(T_x\) is finite almost surely (that is, \(\mathbb{P}(T_x<+\infty)=1\)).
\item[\textnormal{2.}]
Find \(\mathbb{P}(T_3<T_{-1})\). You can use a picture as justification.
\end{examsubparts}
\end{examproblems}
\end{document}
